\section{Introduction}

Image classification is usually formulated as learning one target label from one
image. Natural visual experience is more compositional: multiple objects often
appear in the same field of view, and visual categories are often understood
through comparisons among related sources of evidence. This suggests a
complementary augmentation principle: construct training examples that expose
several labeled sources jointly, while preserving explicit control over how much
visual evidence each source contributes.

Mixed-sample data augmentation moves in this direction by composing examples
and mixing their labels. Mixup~\citep{mixup} and CutMix~\citep{cutmix} are
widely used in modern image-classification recipes~\citep{
wightmanResNetStrikesBack2021,pmlr-v139-touvron21a,
touvronDeiTIIIRevenge2022,liOpenMixupOpenMixup2024}, but they are primarily
pairwise: Mixup globally interpolates two images, and CutMix replaces one
region with a patch from another image. Multi-image methods such as
RICAP~\citep{takahashiDataAugmentationUsing2020} and Mosaic-style
augmentation~\citep{bochkovskiyYOLOv4OptimalSpeed2020} show the value of
composing several images, but typically use a fixed small number of sources and
restricted layout families. Moreover, existing methods usually treat region
area only as a soft-label weight, rather than as a controllable object that can
shape both the visual composition and the supervision signal. Figure~\ref{fig:aug-showcase-clean}
shows a qualitative comparison of these augmentation patterns on matched source
images.

\begin{figure*}[t]
  \centering
  \setlength{\tabcolsep}{2pt}
  \begin{tabular}{ccccc}
    \includegraphics{../evaluation/data/processed/figures/augmentation_showcase/aug_comparison_clean/none__panel1_1024.pdf} &
    \includegraphics{../evaluation/data/processed/figures/augmentation_showcase/aug_comparison_clean/mixup__panel1_1024.pdf} &
    \includegraphics{../evaluation/data/processed/figures/augmentation_showcase/aug_comparison_clean/cutmix__panel1_1024.pdf} &
    \includegraphics{../evaluation/data/processed/figures/augmentation_showcase/aug_comparison_clean/mosaic__panel1_1024.pdf} &
    \includegraphics{../evaluation/data/processed/figures/augmentation_showcase/aug_comparison_clean/treemapmix__panel1_1024.pdf} \\
    (a) None & (b) Mixup & (c) CutMix & (d) RICAP & (e) TreemapMix
  \end{tabular}
  \caption{Qualitative comparison of five augmentation strategies on matched source images using resize-and-center-crop preprocessing only. Mixup blends full images, CutMix pastes a random region, RICAP composes four tiles, and TreemapMix composes $K$ rescaled source images within a randomized layout.}
  \label{fig:aug-showcase-clean}
\end{figure*}


We propose \emph{TreemapMix}, a $K$-way image-mixing augmentation that composes
multiple single-label images into a shared canvas while preserving the training
sample count. TreemapMix samples region weights from a Dirichlet distribution
and maps them to non-overlapping area-proportional regions using a squarified
treemap layout~\citep{squarifiedtreemap}. The Dirichlet concentration
parameter $\alpha$ controls the composition, ranging from dominant-region
mixtures to more balanced multi-image mixtures. To prevent persistent
label-to-area or label-to-rank bias, TreemapMix generates $K$ cyclic views from
each group of $K$ source images, so that every source image is paired with every
sampled region weight across the cyclic orbit. Area-preserving symmetries then
increase spatial diversity without changing the target weights.

TreemapMix also exposes a supervision choice. Region area provides a natural
but imperfect proxy for visual evidence: larger regions usually contain more
pixels, more context, and more object parts, and are therefore often easier to
recognize than smaller regions. This suggests an ordinal learning signal within
a composed image, where labels assigned to larger regions should generally
receive higher scores than labels assigned to smaller regions. However, area is
not a calibrated measure of semantic evidence; a large region may contain mostly
background, while a small region may contain a discriminative object part. We
therefore study two complementary objectives: TreemapMix-SCE, which treats
region areas as soft probability targets and optimizes soft cross-entropy, and
TreemapMix-PL, which converts region areas into area-induced ordinal
supervision by ranking the target classes and optimizing a Plackett-Luce
listwise objective~\citep{luce1959,plackett1975,listmle}.

We evaluate TreemapMix on ImageNet-1K~\citep{
dengImageNetLargescaleHierarchical2009} under matched training recipes across
multiple modern image-classification backbones. Compared with 13 baseline configurations, TreemapMix-PL achieves
the strongest Top-1 accuracy across the evaluated backbones, while
TreemapMix-SCE gives the lowest expected calibration error (ECE). Further
diagnostics show that TreemapMix-PL improves alignment between logits and the
area-induced ordering of composed regions, and COCO transfer experiments
indicate that TreemapMix-pretrained representations remain useful for detection
and instance segmentation~\citep{linMicrosoftCOCOCommon2015}.

Our contributions are:
\begin{enumerate}
    \renewcommand{\labelenumi}{\roman{enumi}.}
    \item We introduce TreemapMix, a Dirichlet-controlled $K$-way image-mixing
    augmentation that supports variable-source composition with explicit
    area-proportional treemap layouts while preserving the training sample
    count.

    \item We formulate two supervision objectives from the same composed image:
    area-proportional soft cross-entropy for probability matching and
    Plackett-Luce ranking for area-induced ordinal supervision.

    \item We analyze how the Dirichlet concentration parameter controls target
    variability, probability-matching supervision, and ranking emphasis.

    \item We provide a matched empirical evaluation on ImageNet-1K, including
    accuracy, calibration metrics, area-logit alignment analysis, and COCO
    detection and instance-segmentation transfer experiments.
\end{enumerate}
